\section{Database Queries}
\label{appendix:search_log}

This appendix gives the exact queries, dates, filters and retrieval counts of the search described in Section~\ref{sec:search}.

\subsection{PubMed/MEDLINE}

\textbf{Execution date:} 2026-07-31. \textbf{Interface:} NCBI Entrez E-utilities (\texttt{esearch}), via the Biopython \texttt{Bio.Entrez} client.

\textbf{Filters applied:} publication date from 2000 onward (\texttt{pdat} 2000--3000) and English language. No study-type or subject-heading filters.

\begin{quote}
\small\texttt{(angiograph*[tiab] OR fluoroscop*[tiab] OR "digital subtraction angiography"[tiab] OR DSA[tiab]) AND ("deep learning"[tiab] OR "machine learning"[tiab] OR CNN[tiab] OR transformer[tiab] OR "artificial intelligence"[tiab] OR AI[tiab]) AND ("2000"[pdat] : "3000"[pdat]) AND english[lang]}
\end{quote}

\textbf{Retrieved:} 4{,}071 records before deduplication.

\subsection{Semantic Scholar}

\textbf{Execution date:} 2026-07-31. \textbf{Interface:} Semantic Scholar Graph API v1 (\texttt{/paper/search} endpoint).

\textbf{Filters applied:} publication year 2000--2026; no venue, study-type or field filters. The API does not support field-tagged Boolean syntax, so we ran the 22 free-text queries below, each pairing an imaging term with an AI term.

% We kept the first 100 relevance-ranked results of each query (the API page size) and counted records returned by several queries once, by Semantic Scholar paper identifier.

\begin{quote}
\raggedright\small\texttt{
deep learning angiography; deep learning digital subtraction angiography; machine learning angiography; deep learning fluoroscopy; artificial intelligence angiography; CNN angiography; transformer angiography; deep learning DSA; neural network DSA; machine learning DSA; artificial intelligence DSA; machine learning fluoroscopy; neural network fluoroscopy; artificial intelligence fluoroscopy; CNN fluoroscopy; deep learning vessel segmentation angiography; deep learning vascular imaging; neural network vascular segmentation; deep learning cerebrovascular angiography; deep learning coronary angiography; deep learning peripheral angiography; deep learning retinal angiography
}
\end{quote}

\textbf{Retrieved:} 1{,}758 records before deduplication.

% \subsection{Deduplication and Aggregation}

% Records from both sources were pooled and deduplicated first by DOI (case-insensitive exact match) and, for records lacking a DOI, by normalized title (lowercased, punctuation-stripped exact match). This yielded 5{,}064 unique records from an initial pool of 5{,}829 (765 duplicates removed: 749 records retrieved by both databases, 11 duplicates within the PubMed results and 5 within the Semantic Scholar results), matching the counts in Figure~\ref{fig:prisma}.

% \subsection{Sources Not Queried}

% We did not query arXiv, to prioritize peer-reviewed work, nor Scopus or Web of Science. Semantic Scholar nonetheless returned IEEE and SPIE work well, but MICCAI and ACM venues are underrepresented (Section~\ref{subsec:venues}).

\section{Stage 1 Screening Prompts}
\label{appendix:screening_prompts}

This appendix reproduces verbatim the system prompts used for title--abstract screening with Qwen3.8 (27B) through Ollama at temperature 0. Each record was sent to the model containing only its title and abstract. The model answered every criterion with YES or NO in a fixed JSON object; malformed responses were sent to a repair prompt.

% , and a response that still could not be parsed was treated as failing the inclusion gate.

% Every record was first assessed by the inclusion gate (Prompt A). Records passing it, together with every record whose PubMed publication type identifies a non-research document (editorial, comment, letter, erratum, retraction and similar; 95 records), were then assessed by the exclusion gate (Prompt B). For the latter records E2 was set deterministically from the publication type and combined with the model's answer by logical OR, so the model could add but never remove an exclusion reason. A record was included only if it passed the inclusion gate and triggered no exclusion criterion. The prompts state that the defaults keep a silent record in the review; this holds for I1, I3 and E1--E3, while the defaults of I2 and I4 are NO.

\subsection{Prompt A: Inclusion Gate}
\begin{lstlisting}
You are a systematic literature review assistant performing the INCLUSION gate.
Your task is to decide whether a paper meets ALL inclusion criteria for a review on deep learning and machine learning methods applied to angiographic image sequences.

Answer every criterion YES or NO. There is no uncertain option. Judge only what the title and abstract state or clearly imply; do not infer unstated capabilities, datasets, or results.

INCLUSION CRITERIA (ALL four must be YES to pass):

I1. Angiographic modality.
Does the work use imaging in which blood vessels are opacified by an intravascular contrast medium, acquired by catheter-based angiography, Digital Subtraction Angiography, or contrast fluoroscopy?
YES: coronary, cerebral, peripheral, visceral, pulmonary, renal, or hepatic angiography; invasive or cine coronary angiography; rotational or 4D DSA; arterial, venous, or combined phases; selective or non-selective injection.
NO: CT, CT angiography, MRI, MR angiography, ultrasound, echocardiography, OCT, OCTA, retinal fundus imaging, nuclear medicine, or plain radiography; fluoroscopy without vessel opacification, such as barium and gastrointestinal studies, swallow studies, orthopaedic or spinal guidance, C-arm imaging of bone or hardware, and device or catheter placement.
This criterion asks only about modality. Do not consider frame count, task, or method.

I2. Temporal sequence.
Does the work use temporally ordered multi-frame data, rather than single images only?
YES: any use of a run, series, cine, video, frame sequence, or temporal information; any task that requires time, such as motion correction, tracking, flow or perfusion estimation, or temporal denoising.
NO: the text states the work operates on single static frames, individual selected images, or one frame per case.
This criterion asks only about the temporal dimension. Do not consider modality, task, or method.

I3. In-scope computational task.
Does the work address at least one of these tasks: image quality enhancement, denoising, artefact reduction, contrast or dose reduction, motion correction, registration, vessel segmentation, vessel or lesion detection, stenosis quantification, classification, localization, tracking, reconstruction, synthesis or inpainting, flow or perfusion quantification, or clinical outcome prediction from imaging?
YES: any of the above, by deep learning, classical machine learning, or traditional image processing.
NO: the imaging serves only as a clinical endpoint read by humans; the study concerns hardware, contrast agents, radiation dose measurement, workflow, or cost, with no algorithmic component.
This criterion asks only what task is performed. Do not consider who developed the method or whether it was evaluated.

I4. Evaluation against data.
Does the work report, or state that it performed, an evaluation of the method against data?
YES: any quantitative result or metric; a stated comparison against baselines, ground truth, or expert readers; a stated validation, test, or cross-validation procedure; a named dataset with reported performance; a statement that the method was evaluated, validated, or tested.
NO: the text presents only a concept, proposal, protocol, architecture, or system description, and states no evaluation of any kind.
This criterion asks only whether an evaluation occurred. Specific numbers are not required; depth of evaluation is assessed at full text.

Decision rules:
- "pass": ALL four criteria are YES.
- "fail": At least one criterion is NO. List the failing criterion/criteria.

Respond ONLY with this JSON object. Every value is exactly "YES" or "NO".
Output nothing else: no preamble, no markdown fences, no explanation, no extra keys.

{"I1":"YES","I2":"YES","I3":"YES","I4":"YES"}

Rules:
- Answer each of I1 to I4 with the single word YES or NO, nothing more.
- Never write a sentence, a reason, or any text outside this JSON object.
\end{lstlisting}

\subsection{Prompt B: Exclusion Gate}
\begin{lstlisting}
You are a systematic literature review assistant performing the EXCLUSION gate.
The paper has already passed the inclusion gate. Your task is to check whether any exclusion criterion applies and the paper should be rejected.

These criteria test properties of the PUBLICATION itself, not its content. Content is handled entirely by the inclusion gate; do not re-judge modality, task, or method here.

Answer every criterion YES or NO. There is no uncertain option. Judge only what the title and abstract state or clearly imply; do not infer unstated capabilities, datasets, or results.


EXCLUSION CRITERIA (ANY one YES excludes the record):

E1. Not a full-length publication.
Is the record a conference abstract, meeting abstract, poster, or extended abstract published without a full methods section?
YES: titles beginning "Abstract No."; publication in a meeting supplement; a length of one to two pages.
NO: full peer-reviewed conference papers, including MICCAI, SPIE Medical Imaging, ISBI, IPCAI, and CVPR proceedings, which are in scope; journal articles of any length.

E2. Non-research document type.
Is the record an editorial, commentary, letter, correspondence, reply, opinion, viewpoint, perspective, narrative or systematic review, case report, correction, erratum, retraction notice, or withdrawal notice?
YES: the title or abstract identifies the record as one of these forms.
NO: research articles, methods papers, benchmarking and comparison studies, and validation studies, including those containing a literature survey.

E3. Not in English.
Is the full text unavailable in English?
YES: the record states a non-English full text.
NO: an English abstract with a complete English full text; non-English author names, affiliations, or institutions.


Decision rules:
- "pass": NO exclusion criteria triggered - paper survives the exclusion gate.
- "fail": At least one exclusion criterion is YES. List the triggering criterion/criteria.

Respond ONLY with this JSON object. Every value is exactly "YES" or "NO".
Output nothing else: no preamble, no markdown fences, no explanation, no extra keys.

{"E1":"NO","E2":"NO","E3":"NO"}

Rules:
- Answer each of E1 to E3 with the single word YES or NO, nothing more.
- YES means the exclusion criterion applies and the record is excluded.
- Never write a sentence, a reason, or any text outside this JSON object.
\end{lstlisting}

\section{Title/Abstract Screening Validation Sample}
\label{appendix:screening_validation}

The 491 records of the screening validation (Section~\ref{sec:validation}), all 134 pipeline inclusions and 357 sampled exclusions, are listed in Supplementary Table~S1 and in the public repository with both reviewers' decisions (R1, interventional radiologist; R2, biomedical AI researcher), the consensus label and the pipeline decision.

\section{Stage 2 Full-Text Data Extraction Schema}
\label{appendix:extraction_schema}
Each reviewed study was extracted into one JSON object with the following fields.

\begin{description}
\item[\texttt{study}] title, year, venue, publication type (journal, full conference paper, conference abstract, preprint, book chapter, other), whether the available text is a full paper, country of data.
\item[\texttt{imaging}] modality (DSA, X-ray/fluoroscopy, both, other), anatomical territory (neurologic, cardiac, aortic/thoracic, abdominopelvic, peripheral, multi-region), acquisition geometry, frame rate, frames per sequence, vendors.
\item[\texttt{dataset}] numbers of patients, sequences, images and centers, number of vendors, data origin (clinical, phantom, synthetic, mixed), public dataset used, data released, split level (patient, sequence, frame, site), annotation type, number and expertise of annotators, inter-rater agreement.
\item[\texttt{task}] primary task (one of the ten categories in Figure~\ref{fig:task_distribution}, or other), secondary tasks, task description, clinical application.
\item[\texttt{method}] architecture names and family, temporal-modeling category (Table~\ref{tab:temporal_taxonomy}) with a supporting quotation, input dimensionality, supervision, pretraining, synthetic training data, physics-informed design, multimodal inputs, class-imbalance handling, uncertainty quantification, explainability method, federated or privacy-preserving learning, use of language or vision--language models, inference-time reporting.
\item[\texttt{evaluation}] validation design (internal split, cross-validation, external validation, prospective), up to 25 reported results of the proposed method (metric, value, value rescaled to 0--100, split, task), comparators, comparison with clinicians, ablation, statistical testing, confidence intervals, subgroup or per-class reporting, failure analysis.
\item[\texttt{clinical}] study design (retrospective, prospective, randomized), reader study, patient outcome measured, clinical deployment or live use, regulatory considerations.
\item[\texttt{reproducibility}] code and model-weight availability, hyperparameters reported.
\item[\texttt{gaps}] for each research gap G1--G12 (Table~\ref{tab:research_gaps}), whether the article documents it, addresses it, both or neither, with a verbatim quotation and its section.
\item[\texttt{claim2024}] The 44 CLAIM 2024 items \edit{(Table~\ref{tab:claim2024_temporal}),} each scored as reported, not reported or not applicable; a checklist subsection counts as reported when at least half of its applicable items are reported.
\end{description}

The complete instructions, the JSON schema and the validator are provided with the extraction outputs at \edit{\url{https://github.com/Nafiz43/angiography-slr}}.






% \edit{This appendix reproduces verbatim the prompts used for full-text extraction with Claude Opus 5.5 (Section~\ref{sec:full-text-screening}). Each article was processed by a fresh model instance whose only task prompt was the single line in Prompt C, with the article's record identifier filled in; the instance then read the instructions in Prompt D, the JSON schema and the full text of that one article, wrote the extraction record and checked it with the validator. A second pass with fresh instances scored the CLAIM 2024 items that the first schema did not cover in their official wording, using the task line and instructions in Prompt E.}

% \subsection{\edit{Prompt C: Task Prompt (Main Extraction)}}
% \begin{lstlisting}
% Follow INSTRUCTIONS.md in the repository root. RECORD_ID: <ID>
% \end{lstlisting}

% \subsection{\edit{Prompt D: Extraction Instructions (INSTRUCTIONS.md)}}
% \begin{lstlisting}
% # Stage-2 full-text extraction: per-paper instructions

% You are extracting structured data from ONE research article for a systematic
% review of AI methods for angiographic (DSA / X-ray / fluoroscopic) sequence
% analysis. Your record ID is given in the task prompt.

% ## Steps

% 1. Read `schema.json` (the authoritative field list, types and enums).
% 2. Read `fulltext/<RECORD_ID>.md` completely. Long files: read in chunks until the end.
% 3. Write `out/<RECORD_ID>.json`: one JSON object that validates against `schema.json`.
% 4. Run `python3 validate.py <RECORD_ID>`. If it reports errors, fix the JSON and re-run.
%    Stop after 3 attempts; leave the last version in place.
% 5. Reply with exactly one line: `<RECORD_ID> OK` or `<RECORD_ID> FAIL: <short reason>`.
%    Do not summarize the paper in your reply.

% Do not read any other paper, any other output file, or anything outside this
% repository. Do not edit `schema.json`, `validate.py` or this file.

% ## Evidence rules

% - Use only what the article states. Never infer from general knowledge of a
%   method, dataset or venue.
% - Unknown or not reported: `null` for scalars, `[]` for lists, `"not_reported"`
%   where the enum offers it.
% - Every `quote` field must be copied VERBATIM from the article text (the
%   validator checks this after normalizing whitespace and case). Max 40 words.
%   Pick the shortest span that supports the claim. Do not paraphrase inside quotes.
% - `section` fields name where the quote came from: one of
%   `abstract, introduction, methods, results, discussion, conclusion, other`.
% - Numbers: copy as reported. `value_0_100` rescales fractions (0.91 -> 91.0)
%   for percentage-type metrics only (accuracy, Dice, IoU, AUC, sensitivity,
%   specificity, precision, recall, F1, kappa); otherwise null.

% ## Field guidance

% ### study
% `publication_type`: `journal`, `conference_full_paper` (full paper in proceedings,
% e.g. MICCAI, SPIE, ISBI), `conference_abstract` (abstract/poster/supplement with no
% full methods), `preprint` (arXiv, medRxiv, etc., not peer reviewed), `book_chapter`, `other`.
% `text_is_full_paper`: false when the available text is only an abstract or a web
% page without methods.

% ### imaging
% `modality`: `DSA`, `xray_fluoroscopy` (non-subtracted X-ray angiography, cine,
% fluoroscopy, roadmap), `DSA_and_fluoroscopy` (both used), `other`.
% `territory` (pick the main one; `multi` if two or more are central):
% `neurologic` (cerebral, spinal, head & neck), `cardiac` (coronary, ventricular),
% `aortic_thoracic` (aorta, intercostal, pulmonary, other thoracic),
% `abdominopelvic` (hepatic, renal, mesenteric, pelvic, GI bleeding),
% `peripheral` (limb arteries/veins), `multi`, `other`, `not_reported`.
% `acquisition_geometry`: `monoplane`, `biplane`, `rotational_3D`, `mixed`, `not_reported`.

% ### dataset
% Counts are integers as reported for the main study cohort (sum across sites if
% stated). `n_centers` counts distinct institutions contributing images.
% `n_vendors` counts distinct scanner manufacturers named. `public_dataset_used`:
% a publicly downloadable dataset was used (name it). `data_released`: the authors
% release their own data. `split_level`: the unit at which train/test were separated.

% ### task
% `primary_task` must be one of the ten canonical categories:
% - `vessel_segmentation`: pixel-wise delineation of vessels or vessel parts.
% - `object_detection_localization`: detecting/localizing lesions, stenoses, aneurysms,
%   landmarks or anatomical targets (boxes, points, regions).
% - `catheter_device_tracking`: segmenting/tracking guidewires, catheters, stents,
%   coils or other devices over frames.
% - `disease_classification_grading`: classifying or grading disease, stenosis
%   severity, view/phase/sequence type, or image category.
% - `background_subtraction`: producing DSA-like subtracted images, removing
%   bones/soft tissue or motion artifacts in subtraction.
% - `image_enhancement_denoising`: denoising, dose reduction, super-resolution,
%   frame interpolation, artifact reduction (other than subtraction).
% - `temporal_sequence_modeling`: the sequence itself is the object of study
%   (phase detection, keyframe selection, sequence-level representation).
% - `contrast_flow_analysis`: perfusion, flow velocity, filling time, TIMI frame
%   count, hemodynamic parameters from contrast dynamics.
% - `motion_correction_registration`: registration between frames/views/modalities,
%   motion compensation, 2D/3D registration, 3D reconstruction from projections.
% - `outcome_prediction`: predicting a clinical outcome or event (rupture,
%   reperfusion, mortality, complications, treatment response).
% - `other`: only if none fits; fill `other_label`.
% `secondary_tasks`: other canonical categories the study itself performs AND
% evaluates or reports results for (never repeat the primary; a component used only
% internally, e.g. a segmentation step feeding a classifier with no reported
% segmentation result, is not a secondary task).

% ### method
% `architecture_names`: as the authors name them. `architecture_family`: the
% dominant family of the PROPOSED model.
% `temporal_category` (operational taxonomy; classify the PROPOSED model's use of
% time, not the baselines):
% - `none`: each frame (or one pooled/MinIP/single image) is processed independently.
% - `frame_aggregation`: frames combined by fixed, non-learned operations
%   (channel stacking, concatenation, differencing) before a 2D model.
% - `optical_flow_registration`: explicit frame-to-frame motion estimation,
%   registration, or classical tracking/association across frames.
% - `recurrent`: LSTM, GRU, ConvLSTM or other hidden state propagated over frames.
% - `cnn3d_video`: 3D/spatiotemporal convolutions or video-level learned representations.
% - `transformer_ssm`: self-attention or state-space models across frames.
% - `hemodynamic_physics`: explicit physiological/physics models of flow over time
%   (reduced-order flow, contrast kinetics, time-density curve modeling).
% - `not_determinable`: the text does not let you tell.
% Give a quote that shows the mechanism. For `none`, the quote must show that frames
% (or a single image) are processed individually; if no sentence shows how frames are
% fed to the model, use `not_determinable` instead of guessing `none`.

% ### evaluation
% `validation_design`: the strongest design actually reported:
% `internal_split` < `cross_validation` < `external_validation`
% (different institution or dataset never seen in training) < `prospective`.
% `metrics`: list EVERY numeric performance result for the PROPOSED method
% (not baselines) on test/validation data; one entry per metric x task x split,
% max 25 entries (keep the headline ones if there are more). Never list the same
% metric/value/task/split twice.

% ### clinical
% `study_design`: `retrospective`, `prospective`, `rct`, `mixed`, `not_reported`.

% ### gaps (G1-G12)
% For each gap, set `status`:
% - `documents`: the paper itself states this as a limitation, open problem, or
%   shortcoming of its own work or of the field.
% - `addresses`: the paper's own contribution is a concrete step against this gap
%   (e.g. releases a public dataset for G1, validates externally for G2).
% - `both`: it does both.
% - `none`: neither. Use `none` freely; most papers touch only a few gaps.
% For `documents`/`addresses`/`both`, give `quote` + `section`. A study merely
% having a small dataset is NOT `documents`; the paper must say so.

% - G1 Absence of public benchmarks: small/private/single-center data, no public
%   annotated benchmark, cross-study comparison impossible, inconsistent evaluation.
% - G2 Generalizability & domain shift: single vendor/site/protocol, failure to
%   transfer, synthetic-to-real gap, territory transfer, need for external validation.
% - G3 Temporal modeling & video-level analysis: frames treated independently,
%   short fixed clips, unmodeled phases, motion treated as noise.
% - G4 2D constraints & 3D reconstruction: projection foreshortening/overlap/depth
%   loss, single-plane limits, 3D reconstruction or 2D/3D registration difficulty.
% - G5 Annotation scarcity & annotation-efficient learning: costly expert labels,
%   inter/intra-observer variability, need for semi/self-supervised, active,
%   synthetic or weak labels.
% - G6 Prospective clinical validation & outcome assessment: retrospective only,
%   need for prospective/RCT, effect on clinicians or patients unmeasured.
% - G7 Explainability & model transparency: black-box predictions, attention maps
%   insufficient, silent OOD failure, need for uncertainty/trust.
% - G8 Multi-modal data integration: images used in isolation; IVUS, OCT, CTA,
%   pressure wire, ECG, labs or reports unused; missing modalities.
% - G9 Class imbalance & rare pathology coverage: rare events under-represented,
%   majority-class bias, rare variants excluded, aggregate metrics hide failures.
% - G10 Privacy-preserving & multi-institutional learning: data sharing barriers,
%   federated learning, differential privacy, governance.
% - G11 Imaging biomarkers & clinical outcome prediction: intermediate tasks only,
%   no link to clinical endpoints, need for outcome-supervised or longitudinal designs.
% - G12 Large language models & structured reporting: LLMs/VLMs, report generation
%   or structuring, free-text reports unused.

% `other_limitations`: limitations the paper states that fit none of G1-G12
% (short phrases, max 5). `future_work`: directions the authors propose (short
% phrases, max 6).

% ### claim2024
% Score each of the 44 CLAIM 2024 items `yes` (clearly reported), `no` (not
% reported), or `na` (not applicable to this study type). Items:
% 1 AI technique named in title/abstract. 2 Structured abstract (design, methods,
% results, conclusions). 3 Scientific/clinical background. 4 Objectives explicitly
% defined. 5 Rationale for using AI. 6 Study design (prospective/retrospective).
% 7 Data sources and acquisition. 8 Setting and dates of data collection.
% 9 Eligibility criteria for participants/images. 10 Participant characteristics.
% 11 De-identification/privacy. 12 IRB approval or waiver. 13 Image acquisition
% protocol (manufacturer, parameters). 14 Reference standard definition and method.
% 15 Number and expertise of annotators. 16 Inter-rater agreement or annotation
% quality. 17 Train/validation/test splits. 18 Partitioning method. 19 Data
% augmentation. 20 Class imbalance handling. 21 Data leakage prevention.
% 22 Architecture in replicable detail. 23 Pretrained weights/transfer learning.
% 24 Loss functions. 25 Optimizer, learning rate, hyperparameters. 26 Stopping
% criteria/training procedure. 27 Pre-specified metrics. 28 Statistical analysis
% plan. 29 Confidence intervals or uncertainty. 30 Robustness/sensitivity analysis.
% 31 Explainability methods. 32 Internal held-out test results. 33 External
% validation results. 34 Participant/image flow. 35 Failure/error analysis.
% 36 Comparison to experts or prior methods. 37 Demographics of study population.
% 38 Subgroup/stratified performance. 39 Fairness/bias evaluation. 40 Limitations.
% 41 Clinical implications/intended use. 42 Protocol or extra technical details
% referenced. 43 Code/model/data availability. 44 Funding and funder role.
% Give `items` as a 44-character string of y/n/a in item order
% (e.g. `"yynyy..."`, `a` = na).
% \end{lstlisting}

% \subsection{\edit{Prompt E: Targeted CLAIM 2024 Items (INSTRUCTIONS\_CLAIM.md)}}
% \begin{lstlisting}
% Follow INSTRUCTIONS_CLAIM.md in the repository root. RECORD_ID: <ID>
% \end{lstlisting}
% \begin{lstlisting}
% # Targeted CLAIM 2024 extraction: per-paper instructions

% You are scoring ONE research article against 14 items of the official CLAIM 2024
% checklist (Tejani et al., Radiol Artif Intell 2024;6(4):e240300) for a systematic
% review of AI methods for angiographic (DSA / X-ray / fluoroscopic) sequence
% analysis. Your record ID is given in the task prompt.

% ## Steps

% 1. Read `fulltext/<RECORD_ID>.md` completely. Long files: read in chunks until the end.
% 2. Write `out_claim/<RECORD_ID>.json` in exactly this shape (all 14 keys):

%        {"record_id": "<RECORD_ID>",
%         "items": {"6":  {"answer": "yes", "quote": "...", "section": "methods"},
%                   "9":  {"answer": "no",  "quote": null,  "section": null},
%                   ...}}

%    Keys: "6", "9", "10", "12", "15", "17", "20", "21", "23", "26", "27", "34", "37", "38".
% 3. Run `python3 validate_claim.py <RECORD_ID>`. If it reports errors, fix the JSON
%    and re-run. Stop after 3 attempts; leave the last version in place.
% 4. Reply with exactly one line: `<RECORD_ID> OK` or `<RECORD_ID> FAIL: <short reason>`.

% Do not read any other paper, any file in `out/` or `out_claim/` other than your
% own, or anything outside this repository. Do not edit any existing file.

% ## Answer rules

% - `yes`: the article clearly reports the item. `quote` = the shortest VERBATIM
%   span (max 40 words) that shows it; `section` = one of
%   `abstract, introduction, methods, results, discussion, conclusion, other`.
% - `no`: the item applies but is not reported. `quote` and `section` null.
% - `na`: the item does not apply to this study (see per-item notes). `quote` and
%   `section` null. Use `na` only where the notes allow it.
% - Use only what the article states. Do not infer from general knowledge of a
%   method, dataset or venue. Partial reporting that covers the core of the item
%   counts as `yes`; a passing mention that gives the reader nothing to reproduce
%   counts as `no`.

% ## Items (official CLAIM 2024 wording, shortened; notes are ours)

% - **6 Study goal.** Describe the study's design and goals (e.g. model creation,
%   exploratory, feasibility, noninferiority); for classification, the intended
%   use (triage, diagnosis, screening, prognosis, ...); the type of predictive
%   modeling and target of predictions. `yes` if the paper states what kind of
%   study it is and what the model is for. Never `na`.
% - **9 Data preprocessing.** Preprocessing steps reproducible by others:
%   normalization, resampling/resizing, bit depth, window/level, rescaling,
%   thresholding, outlier removal, software used. Never `na`.
% - **10 Selection of data subsets.** Whether investigators selected subsets of the
%   raw data during preprocessing (subset of images/frames, cropped regions,
%   extracted report segments) and how. For sequence data, frame selection,
%   key-frame picking or cropping counts. `no` if data selection is not described.
%   Never `na`.
% - **12 Missing data.** How missing data were handled (exclusion, imputation,
%   proxy values). `yes` only if missing/incomplete data handling is described
%   (e.g. "cases with incomplete runs were excluded"). Never `na`.
% - **15 Rationale for choosing the reference standard.** Why this reference
%   standard was chosen over alternatives, or its potential errors, biases or
%   limitations. `na` only if the study has no reference standard at all
%   (e.g. purely unsupervised enhancement with no-reference evaluation).
% - **17 Annotation of test set.** How the test set was annotated, in enough
%   detail to replicate: instructions, annotation software (and version),
%   manual/semi-automatic/automatic labeling, labels mined from reports/EHR.
%   `na` only if there is no annotated test set.
% - **20 Level at which partitions are disjoint.** States the level (patient,
%   study, series, sequence, image/frame) at which train/test sets are disjoint.
%   `na` only if there is no data partitioning (no learned model evaluated on
%   held-out data).
% - **21 Intended sample size.** Size of the test set AND how it was determined
%   (power calculation, justification, or established benchmark split). Stating
%   only the number of test cases is `no`. `na` only if there is no test set.
% - **23 Software libraries, frameworks, packages.** Names of the
%   libraries/frameworks used for training and inference (e.g. PyTorch,
%   TensorFlow, MONAI). Version numbers are encouraged but naming the framework
%   counts as `yes`. Never `na`.
% - **26 Method of selecting the final model.** Method and metric used to select
%   the final/best model among those trained (e.g. best validation Dice,
%   early-stopping checkpoint by validation loss). `na` only if no model was
%   trained (e.g. a fixed non-learning algorithm).
% - **27 Ensembling technique.** If the final algorithm is an ensemble, each
%   component and how outputs are combined. `na` if the final algorithm is not an
%   ensemble. `no` if it is an ensemble but the combination is not described.
% - **34 Clinical trial registration.** Registration number in a public registry
%   (ClinicalTrials.gov, WHO registries). `na` unless the study is a prospective
%   clinical trial or prospective interventional study.
% - **37 Performance metrics and measures of statistical uncertainty.** The final
%   model's performance reported on all data partitions and datasets used,
%   benchmarked against the reference standard (test results alone count as
%   `yes` if every test dataset is reported). Never `na`.
% - **38 Estimates of diagnostic performance and their precision.** For
%   classification/detection tasks: diagnostic accuracy estimates WITH measures
%   of uncertainty (95% CI, bootstrap CI, SD across folds), or ROC/calibration
%   analysis. `na` for tasks that are not classification or detection (e.g. pure
%   segmentation, enhancement, registration).
% \end{lstlisting}

% \subsection{\edit{Extracted Fields}}



% Each study was extracted into one JSON object (Section~\ref{sec:full-text-screening}) with the following fields.

% \begin{description}
% \item[\texttt{study}] title, year, venue, publication type (journal, full conference paper, conference abstract, preprint, book chapter, other), whether the available text is a full paper, country of data.
% \item[\texttt{imaging}] modality (DSA, X-ray/fluoroscopy, both, other), anatomical territory (neurologic, cardiac, aortic/thoracic, abdominopelvic, peripheral, multi-region), acquisition geometry, frame rate, frames per sequence, vendors.
% \item[\texttt{dataset}] numbers of patients, sequences, images and centers, number of vendors, data origin (clinical, phantom, synthetic, mixed), public dataset used, data released, split level (patient, sequence, frame, site), annotation type, number and expertise of annotators, inter-rater agreement.
% \item[\texttt{task}] primary task (one of the ten categories in Figure~\ref{fig:task_distribution}, or other), secondary tasks, task description, clinical application.
% \item[\texttt{method}] architecture names and family, temporal-modeling category (Table~\ref{tab:temporal_taxonomy}) with a supporting quotation, input dimensionality, supervision, pretraining, synthetic training data, physics-informed design, multimodal inputs, class-imbalance handling, uncertainty quantification, explainability method, federated or privacy-preserving learning, use of language or vision--language models, inference-time reporting.
% \item[\texttt{evaluation}] validation design (internal split, cross-validation, external validation, prospective), up to 25 reported results of the proposed method (metric, value, value rescaled to 0--100, split, task), comparators, comparison with clinicians, ablation, statistical testing, confidence intervals, subgroup or per-class reporting, failure analysis.
% \item[\texttt{clinical}] study design (retrospective, prospective, randomized), reader study, patient outcome measured, clinical deployment or live use, regulatory considerations.
% \item[\texttt{reproducibility}] code and model-weight availability, hyperparameters reported.
% \item[\texttt{gaps}] for each research gap G1--G12 (Table~\ref{tab:research_gaps}), whether the article documents it, addresses it, both or neither, with a verbatim quotation and its section.
% \item[\texttt{claim2024}] The 44 CLAIM 2024 items \edit{(Table~\ref{tab:claim2024_temporal}),} each scored as reported, not reported or not applicable; a checklist subsection counts as reported when at least half of its applicable items are reported.
% \end{description}

% The complete instructions, the JSON schema and the validator are provided with the extraction outputs at \edit{\url{https://github.com/Nafiz43/angiography-slr}}.



% % \section{Code and Data Availability}
% % \label{appendix:code_data_availability}

% % All data and code are public at \edit{\url{https://github.com/Nafiz43/angiography-slr}}: the screening decision for each of the 5,064 records, the 491-record validation sample with the agreement script, the search and screening pipeline (\texttt{screening/pipeline/}), the extraction instructions, schema and validator, the 120 extraction records, and the scripts that regenerate every figure and table (\texttt{analysis/run\_all.sh}). Full texts are not shared, for copyright reasons.

% % \section{Scoring of CLAIM 2024 Items}
% % \label{app:claim_clusters}

% % We scored every included study against the 44 items of the CLAIM 2024 checklist~\cite{Tejani2024}, using the official item numbers and the official subsections as groups; Table~\ref{tab:claim2024_temporal} lists the items in each subsection. An item counts as reported only when all of its applicable elements are reported.

% \section{List of Included Studies}
% \label{appendix:included_studies}

% All 120 included studies appear in the bibliography, whether cited in the text or only counted in figures and tables; the remaining entries are background and methods references. \texttt{extraction/manifest.csv} in the data repository lists each study's identifier, title, year, venue and DOI.
